%%%%%%%%%%%%%%%%%%%%%%%%%%%%%%%%%%%%%%%%%%%%%%%%%%%%%%%%%%%%%%%%%%%%%%%%%%%%%%%
%  s00-front.tex —— 封面、题词、摘要、译者说明、凡例、术语表
%%%%%%%%%%%%%%%%%%%%%%%%%%%%%%%%%%%%%%%%%%%%%%%%%%%%%%%%%%%%%%%%%%%%%%%%%%%%%%%

%================================ 封面 =======================================%
\begin{titlepage}
\centering
\vspace*{2.2cm}

{\Huge\bfseries 辛拓扑导论\par}
\vspace{0.9em}
{\Large\itshape A Symplectic Prolegomenon\par}

\vspace{2.2em}
{\large Ivan Smith\par}
\vspace{0.4em}
{\small Centre for Mathematical Sciences, University of Cambridge, CB3 0WB, England\par}
\vspace{0.3em}
{\small\texttt{[邮箱隐去]}\par}

\vspace{2.6em}
\begin{minipage}{0.78\textwidth}
\centering
{\large 中文译注版\par}
\vspace{0.6em}
{\small 依据 arXiv:1401.0269v2（2014 年 9 月修订版）翻译\\
附学术级审校批注、本节重点与图表还原\par}
\end{minipage}

\vspace{2.4em}
\begin{minipage}{0.82\textwidth}
\small
\noindent\textbf{内容提要}\quad 一个辛流形会给出一个三角 $A_\infty$-范畴，即导出 Fukaya 范畴；
它编码了拉格朗日子流形以及由 Floer 上同调所探测的动力学信息。本综述旨在让读者对
Fukaya 范畴是什么、它从何而来、以及辛拓扑学家想用它做什么，获得一些直观的理解。
\end{minipage}

\vfill
{\small 2014 年 1 月（v1）；2014 年 9 月修订（v2）\par}
\vspace{0.4em}
{\small 中译本：2026 年 10 月\par}
\end{titlepage}

%================================ 题词 =======================================%
\thispagestyle{empty}
\vspace*{0.28\textheight}
\begin{epigraph}{\itshape
……你想说的一切都需要一个语境。若给出完整语境，人们会觉得你是个絮叨的老傻瓜；
若不给出语境，人们又会觉得你是个言简意赅的老傻瓜。}{\itshape Julian Barnes，《凝视太阳》（Staring at the Sun）}
\end{epigraph}
% [Tutor Note: 原文题词 “everything you wanted to say required a context...” 是全书基调：
%   作者自嘲式地提醒读者，本综述必然在“讲得太多”与“讲得太少”之间取舍。]

\clearpage

%================================ 译者说明 ===================================%
\section*{译者说明}
\addcontentsline{toc}{section}{译者说明}
\markboth{译者说明}{译者说明}

\paragraph{一、底本.}
本译本的底本为 Ivan Smith 的综述文章 \emph{A Symplectic Prolegomenon}，arXiv:1401.0269v2
（\texttt{math.SG}，2014 年 9 月 7 日修订版，47 页，11 幅插图）。该文是为数不多的、
面向“想用 Fukaya 范畴”的读者的整体性导引之一。

\paragraph{二、关于底本的一处说明.}
委托方原始说明中提到 Edward Witten 的文章；经核对，本 PDF 的实际作者是 Ivan Smith（剑桥大学），
主题为辛拓扑与导出 Fukaya 范畴，属于 \texttt{math.SG} 领域的数学综述，并非 Witten 的物理论文。
Witten 在镜像对称、拓扑弦与量子上同调方面的物理论述与本主题密切相关，但本文本身是纯数学的。
译文据实以底本为准，未作任何“向物理靠拢”的改写。

\paragraph{三、翻译原则.}
\begin{enumerate}
  \item \textbf{数学无损}：全部公式、上下标、指标、$\bZ_2$ 分次、符号约定均按底本原样保留；
        只把叙述性文字译为中文。
  \item \textbf{术语统一}：采用中文学界约定俗成的译名（如 gauge theory 规范场论、holonomy 和乐、
        anomaly 反常、Floer 上同调、Fukaya 范畴），人名与专名保留英文（详见文末术语表）。
  \item \textbf{不作增删}：不添加底本没有的定理、例子、文献或物理图像；对作者的“显然可得”、
        “标准论证”等跳步，只在批注中指出其依赖的核心工具。
  \item \textbf{编号一致}：公式编号按节编号，与底本完全对应（如 (3.7)、\eqref{Eqn:HMS}），
        便于与原文逐式对照。
\end{enumerate}

\paragraph{四、审校批注（Tutor Note）的用法.}
源码中插入了大量整行的 LaTeX 注释，形如
\begin{quote}\ttfamily
\% [Tutor Note: 此处推导极易出现符号错误，请注意核对原文的 $\cdots$ 约定。]
\end{quote}
这些批注\textbf{不会出现在编译后的 PDF 中}，是留给读者的“旁批”，分为三类：
(i) 易错预警（指标升降、张量缩并、分次符号、定向约定）；
(ii) 含糊点拨（作者写 “It is obvious/standard that\dots” 时，指出其核心逻辑与可查阅的文献类型）；
(iii) 概念辨析（如 exact 与 monotone、$HF^*$ 与 $QH^*$、twist functor 与 Dehn twist）。
在源码目录中检索 \texttt{Tutor Note} 即可定位全部批注；随交付附有《审校批注索引》一份。

\paragraph{五、排版说明.}
\begin{itemize}
  \item 底本使用 \texttt{mathabx} 宏包，该宏包只提供 Metafont 字型，XeLaTeX 无法嵌入，
        故本译本略去；正文用到的 \texttt{\textbackslash pitchfork}、\texttt{\textbackslash maltese}
        等符号由 \texttt{amssymb} 提供，含义不变。
  \item 底本第 7 幅图（亏格 2 曲面）原本是 \texttt{picture} 环境的手绘矢量图，
        本译本改用自原 PDF 以 400\,dpi 裁出的图片，图形本身逐点一致。
  \item 其余 11 幅插图同样自原 PDF 精确裁剪，图题译为中文并标注原文位置（Figure~$n$，第 $p$ 页）。
  \item 编译：\texttt{xelatex} 连续两遍（第一遍生成目录与交叉引用）。
\end{itemize}

\paragraph{六、致谢.}
感谢匿名审稿人对原著的改进；本译本的疏漏一概由译者负责。读者若发现公式或术语问题，
欢迎直接在源码上批注修订。

\clearpage

%================================ 凡例 =======================================%
\section*{凡例与符号}
\addcontentsline{toc}{section}{凡例与符号}
\markboth{凡例与符号}{凡例与符号}

\begin{center}
\small
\begin{tabular}{@{}p{0.11\textwidth}p{0.35\textwidth}p{0.12\textwidth}p{0.33\textwidth}@{}}
\toprule
$\bR,\bC,\bZ,\bQ,\bN$ & 实数域、复数域、整数环、有理数域、自然数 & $\bP^n,\ \bK$ & 射影空间、系数域 \\
$\scrF(X)$ & 辛流形 $X$ 的（导出）Fukaya 范畴 & $HF^*(L,L')$ & Floer 上同调群 \\
$CF^*(L,L')$ & Floer 上链复形 & $\mu^k$ & $A_\infty$-乘法（$k$ 个输入、1 个输出） \\
$(X,\omega)$ & 辛流形及其辛形式 & $\tau_V$ & 沿拉格朗日球面 $V$ 的 Dehn 扭转 \\
$D^b(\scrX)$ & 层的（有界）导出范畴 & $D^{\pi}\scrF(X)$ & 分裂闭包的导出 Fukaya 范畴 \\
$QH^*(X)$ & 量子上同调 & $\Tw,\ \Tw^{\pi}$ & 扭曲复形、幂等完备化后的扭曲复形 \\
\bottomrule
\end{tabular}
\end{center}

\noindent 约定：$\bZ_2=\{0,1\}$；分次记号 $[1]$ 表示分次移位；$\mathrm{Don}(X)$ 表示
以（无障碍）拉格朗日子流形为对象的 Donaldson 范畴；无特别说明时，流形均默认光滑、紧致、
可定向，拉格朗日子流形均默认闭且可定向。

\clearpage

%================================ 术语表 =====================================%
\section*{术语表}
\addcontentsline{toc}{section}{术语表（英汉对照）}
\markboth{术语表}{术语表}

\begin{center}
\footnotesize
\begin{longtable}{p{0.285\textwidth}p{0.155\textwidth}p{0.265\textwidth}p{0.155\textwidth}}
\toprule
\textbf{English} & \textbf{中文} & \textbf{English} & \textbf{中文} \\
\midrule
\endfirsthead
\toprule
\textbf{English} & \textbf{中文} & \textbf{English} & \textbf{中文} \\
\midrule
\endhead
\bottomrule
\endfoot
symplectic manifold & 辛流形 & symplectic form & 辛形式 \\
symplectic topology & 辛拓扑 & Lagrangian submanifold & 拉格朗日子流形 \\
Floer cohomology & Floer 上同调 & Floer complex & Floer 复形 \\
Fukaya category & Fukaya 范畴 & derived Fukaya category & 导出 Fukaya 范畴 \\
Donaldson category & Donaldson 范畴 & triangulated category & 三角范畴 \\
$A_\infty$-category & $A_\infty$-范畴 & twisted complex & 扭曲复形 \\
exact triangle & 正合三角 & distinguished triangle & 特异三角 \\
mapping cone & 映射锥 & long exact sequence & 长正合列 \\
holomorphic curve & 全纯曲线 & pseudoholomorphic curve & 伪全纯曲线 \\
almost complex structure & 近复结构 & taming / tame & 驯服 / 驯服的 \\
moduli space & 模空间 & transversality & 横截性 \\
bubbling & 冒泡 & disc bubble & 圆盘泡 \\
Gromov compactness & Gromov 紧性 & monotone & 单调的 \\
exact Lagrangian & 恰当拉格朗日 & Maslov class / index & Maslov 类 / 指标 \\
minimal Chern number & 极小 Chern 数 & Novikov field / ring & Novikov 域 / 环 \\
quantum cohomology & 量子上同调 & Gromov--Witten invariant & Gromov--Witten 不变量 \\
Dehn twist & Dehn 扭转 & vanishing cycle & 消没圈 \\
Lefschetz fibration & Lefschetz 纤维化 & mapping class group & 映射类群 \\
monodromy & 单值（性） & flux group & 通量群 \\
Hamiltonian isotopy & Hamilton 同痕 & symplectomorphism & 辛同胚 \\
holonomy & 和乐 & cotangent bundle & 余切丛 \\
Stein manifold & Stein 流形 & plurisubharmonic & 多重次调和 \\
Lagrangian correspondence & 拉格朗日对应 & Lagrange surgery & 拉格朗日手术 \\
Whitney trick / disc & Whitney 技巧 / 圆盘 & pearly trajectory & 珠串型轨道 \\
obstruction theory & 障碍理论 & unobstructed & 无障碍的 \\
bounding cochain & 界定上链 & virtual perturbation & 虚拟扰动 \\
spectral sequence & 谱序列 & filtration & 滤过 \\
Cardy relation & Cardy 关系 & Yoneda module & Yoneda 模 \\
Hochschild cohomology & Hochschild 上同调 & quasi-isomorphism & 拟同构 \\
split-closed & 分裂闭包的 & split-generate & 分裂生成 \\
strong generator & 强生成元 & generation time & 生成时间 \\
Orlov spectrum & Orlov 谱 & stability condition & 稳定性条件 \\
autoequivalence & 自等价 & coherent sheaf & 凝聚层 \\
matrix factorisation & 矩阵分解 & homological mirror symmetry & 同调镜像对称 \\
Lagrangian fibration & 拉格朗日纤维化 & integral affine manifold & 整仿射流形 \\
toric degeneration & 环簇退化 & anti-canonical divisor & 反典范除子 \\
genus & 亏格 & hyperelliptic curve & 超椭圆曲线 \\
quartic K3 surface & 四次 K3 曲面 & toric Fano variety & 环簇 Fano 簇 \\
moment map & 矩映射 & quadric hypersurface & 二次超曲面 \\
Tate curve & Tate 曲线 & elliptic surface & 椭圆曲面 \\
isotropic submanifold & 迷向子流形 & displaceable & 可移开的 \\
Euler characteristic & 欧拉示性数 & Poincar\'e duality & Poincar\'e 对偶 \\
compactly supported & 紧支的 & convex at infinity & 无穷远处凸的 \\
pseudogroup & 伪群 & contact transformation & 接触变换 \\
Legendre transform & Legendre 变换 & K\"unneth theorem & K\"unneth 定理 \\
quilt theory & quilt 理论 & closed-open string map & 闭-开弦映射 \\
\end{longtable}
\end{center}

\clearpage
